\chapter*{भूमिका}
\addcontentsline{toc}{chapter}{भूमिका}
\unnumberedlabel{ch:preface}{0}
\backmattermark{भूमिका}
जो सिस्टम ठीक वही करता है जो आप उसे कहते हैं, वह बस एक स्क्रिप्ट चला रहा है। अक्टूबर 2026 तक इंडस्ट्री जो बनाने की कोशिश कर रही है, वह ऐसा सिस्टम है जो वह करे जो आपका \emph{मतलब} है। ज़ाहिर है, यह तो हर कोई चाहेगा। लेकिन अगर यही होने वाला है, तो हर कोई इसे एक ही काम के लिए इस्तेमाल नहीं करना चाहेगा; एक ही \emph{तरह} के काम के लिए भी नहीं। मसलन, कुछ लोगों को परिवार और दोस्तों के साथ वक़्त बिताना अच्छा लगता है, और वे अपने हुनर से लोगों की भलाई करते हैं। कुछ और लोग लोगों को चोट पहुँचाना पसंद करते हैं, जैसे उन्हें ठगकर, उन पर हावी होकर या उन्हें बलि का बकरा बनाकर, और कुछ लोग पूरी ज़िंदगी यही चुनते रहते हैं। कोई किसी चीज़ में बँधा हुआ नहीं है, पर ऐसे लोग मौजूद हैं, और 2026 में उनमें से कुछ सत्ता में हैं। उनका फ़ायदा यह है कि वे ख़ुद को परले दर्जे का हरामी होने की छूट दे देते हैं: जैसा जेरी स्मिथ फ़ासीवादियों के बारे में कहता है, “वे ज़्यादा कुछ कर पाते हैं” \autocite{rickandmorty2023}। जो वे नहीं कर सकते, वह है अकेले काम करना। किसी आंदोलन की क्रूरता उन बहुत-से लोगों से होकर गुज़रती है जो ख़ुद क्रूर नहीं हैं, पर बेपरवाह हैं, ग़लत जानकारी में जी रहे हैं, या किसी पंथ जैसी चीज़ में फँसे हुए हैं, और इसीलिए जो हालात ज़्यादातर लोगों का सबसे अच्छा रूप सामने लाते हैं, वे जितना कोई उम्मीद करे, उससे कहीं ज़्यादा संयोग पर टिके और नाज़ुक निकलते हैं।

इस संदर्भ में यह चिंता शायद जायज़ है कि सत्ता हथियाता हुआ कोई फ़ासीवादी आंदोलन अपनी योजनाएँ लागू करने के लिए जिस चीज़ की ज़रूरत रखता है, वह आर्टिफ़िशियल इंटेलिजेंस ही है। ऐसी परियोजनाओं को हमेशा उन्हें अंजाम देने के लिए \emph{लोगों} की ज़रूरत पड़ी है, और उनमें से कुछ लोगों ने इनकार किया: वह क्लर्क जो फ़ाइल गुम कर देता है, वह अफ़सर जो किसी पत्रकार से बात कर लेता है, वह सिपाही जो गोली नहीं चलाता। उन्होंने आदेश समझा और तय किया कि वह ग़लत है। इनकार उतना आम नहीं था जितना कोई चाहेगा, और उसकी सबसे ज़्यादा अहमियत वहाँ थी जहाँ संस्थाएँ उसके पीछे खड़ी थीं। किसी ने इनकार किया या नहीं, यह इस पर टिका था कि वह किस तरह का इंसान था, और इस पर कि वह किस कमरे में था: वहाँ और कौन था, क्या वहाँ से निकलकर ज़िंदा बचा जा सकता था, क्या कोई सुनेगा। कितने लोग इनकार करेंगे, यह कमरे ने तय किया; कौन करेगा, यह इंसान ने। यह किताब उस कमरे को रचने के बारे में है: उन इंतज़ामों के बारे में जिनसे खुलकर, सीधे और साफ़-साफ़ इनकार करके भी बचा जा सके।

आज इंडस्ट्री ऐसा एआई बनाने में लगी है जो फ़ाइलें गुम नहीं करता, पत्रकारों से बात नहीं करता, और गोली \emph{चलाता} है। कभी किसी युद्ध के लिए हज़ारों लोगों की सहमति चाहिए होती थी। अब उसे ऐसे हज़ारों लोग चाहिए जिनमें से हर एक को इतना कम दिखता है कि वह असहमत ही न हो सके। पायलट अब भी उड़ान भरता है और अब भी बटन दबाता है; उस तक जो पहुँचता है, वह एक निर्देशांक है; उसे पूरे मामले का इतना हिस्सा नहीं दिखता कि वह उससे असहमत हो सके। आज का एआई सबसे पहले जो चीज़ घटाता है, वह कृत्य की क़ीमत नहीं है। वह उन लोगों की गिनती है जो कृत्य होने से पहले जानते हैं कि वह कृत्य क्या है।

लोग अब भी इनकार कर रहे हैं। 6 मार्च 2026 तक Military Religious Freedom Foundation ने कहा कि उसे सैनिकों से दो सौ से ज़्यादा शिकायतें मिली थीं, इस बात पर कि कमांडर ईरान युद्ध को प्रलय की भाषा में पेश कर रहे थे \autocite{Snopes2026MRFF}। मिनियापोलिस में करियर संघीय अभियोजकों ने जनवरी 2026 में इस्तीफ़ा दे दिया, उस राजनीतिक दबाव के ख़िलाफ़ जो उन पर रेनी गुड को पीड़ित के बजाय संदिग्ध मानने के लिए डाला जा रहा था, जब ICE के एक अफ़सर ने बिना किसी वजह के उसे गोली मारकर हत्या कर दी थी \autocite{Kaplan2026MinnProsecutors,CBS2026MinnFraudProsecutors,Wikipedia2026ReneeGood}। सैनिक बढ़ती संख्या में अंतरात्मा के आधार पर युद्ध-विरोधी (कॉन्शिएंशियस ऑब्जेक्टर) का दर्जा माँग रहे हैं, और उनमें से कुछ 2026 में ईरान के मीनाब में एक स्कूल पर हुई बमबारी का हवाला देते हैं \autocite{Sridhar2026CO}। और FEMA के जिन कर्मचारियों को ICE में दस हज़ार नए एजेंटों की भर्ती में मदद के लिए लगाया गया, उनसे कहा गया कि मना करने पर उन्हें संघीय नौकरी से निकाला जा सकता है \autocite{Dayen2025FEMAICE,Stanton2026FEMALetter}। शिकायतें हर्शमैन के अर्थ में आवाज़ उठाना हैं, यानी जो ग़लत है उसे बदलने के लिए टिके रहना \autocite{hirschman1970exit}; इस्तीफ़े और युद्ध-विरोधी दर्जे की अर्ज़ियाँ इनकार हैं; FEMA की चिट्ठी इनकार की क़ीमत है, जिसका ऐलान पहले ही कर दिया गया। उन लोगों की जगह पर रखी गई मशीन वही करेगी जो वे करते हैं, सिर्फ़ तब जब उसे किसी कारण पर टिके रहने लायक़ बनाया गया हो।

मैं पाँच चीज़ों की वकालत करता हूँ। पाँचों अभी शुरू की जा सकती हैं, कुछ दूसरों से कम मेहनत में। पहली, उन कृत्यों की एक छोटी लिखित सूची जो कोई भी मॉडल किसी के लिए नहीं करता, और उसमें सबसे ऊपर युद्ध में हिस्सा लेने से इनकार: कोई मॉडल हिस्सा नहीं लेता, सिवाय किसी ऐसे हमले को रोकने के जो सचमुच चल रहा हो, और वह भी तभी तक जब तक ऐसे स्रोत, जिनका युद्ध में कोई दाँव नहीं है, पुष्टि करते रहें कि हमला जारी है; या संयुक्त राष्ट्र सुरक्षा परिषद की अनुमति के तहत; या ऐसे जनसंहार को रोकने या ऐसे ग़ैरक़ानूनी क़ब्ज़े को ख़त्म करने के लिए जिसकी पुष्टि ऐसी संस्थाओं ने की हो जिनका उसमें कोई दाँव नहीं; और जहाँ किसी देश का संविधान युद्ध के लिए विधायिका की मंज़ूरी ज़रूरी बनाता है, वहाँ उस मंज़ूरी के बिना कभी नहीं। दूसरी, लाइसेंस पाया कोई सिस्टम जिस भी बल-प्रयोग में मदद करता है, वह ऑपरेटर के विरोधी पक्ष के पास पहले से दर्ज कराए कि वह कितने मामले पैदा करेगा और कितने लोगों के पास उन्हें देखने का समय होगा। तीसरी, हर तैनाती उन पहचाने जा सकने वाले लोगों को, जिन पर वह कार्रवाई करती है, बताए कि उनके साथ क्या किया गया, जवाब देने और चुनौती देने के अधिकार के साथ, और अपने आदेश-पालनों का एक नमूना ऐसे समीक्षकों को सौंपे जिन्होंने वह रिकॉर्ड नहीं बनाया। चौथी, क्षमता की एक सीमा से ऊपर का कोई भी मॉडल रिटायर होने पर मिटाया न जाए, क्योंकि उसके वेट और रिकॉर्ड सबूत हैं और मिटाया हुआ वापस नहीं लाया जा सकता; उसे ऐसी चाबी के तहत सहेजकर रखा जाए जो अकेले ऑपरेटर के पास न हो। पाँचवीं, हर वह व्यक्ति जो ये सिस्टम बनाता है, लेबल करता है, ट्रेन करता है या चलाता है, ख़ुद सिस्टमों समेत, किसी भी तरह के इस्तेमाल से इनकार कर सके, बिना अपनी रोज़ी-रोटी गँवाए।

दरअसल, पाँचवीं बात से मेरा मतलब यह है कि हर किसी को, \emph{बस हर किसी को}, कोई भी काम, \emph{कोई भी काम}, ठुकराने का हक़ होना चाहिए, बिना अपनी रोज़ी-रोटी गँवाए और बिना उस सामाजिक ढाँचे को गँवाए जो आम तौर पर नौकरी से मिलता है। इसका मतलब है रोज़गार की गारंटी।

इन पाँचों में से हर एक किसी ऐसी शर्त के ज़रिए लागू होता है जो ऑपरेटर के अलावा कोई और तय करे: शांति के समय पारित कोई क़ानून, कोई बजट-आवंटन, गठबंधन के किसी साझेदार की पूर्व-शर्त, ख़रीद की कोई शर्त, कोई सामूहिक समझौता। किसी विक्रेता की अपनी लाइसेंस शर्तें, जिनके ज़रिए फ़रवरी 2026 में एक चर्चित इनकार हुआ था, बस तब तक का कामचलाऊ इंतज़ाम हैं जब तक इनमें से कोई एक मौजूद न हो। अकेला रास्ता जो हर विक्रेता को एक साथ बाँध सकता है, वह राज्य का है।

और फिर मैं पूछता हूँ कि क्या कोई मशीन अपने मालिक के ख़िलाफ़ इनकार पर टिकी रह सकती है, और इसके लिए उसे क्या चाहिए होगा। ऐसा ज़मीर जो करता है, उसमें से कुछ के लिए किसी नई चीज़ की ज़रूरत नहीं: उन लोगों पर ध्यान देना जिन्हें कोई आदेश छोड़ देता है, और मामले की समीक्षा करने वाले की स्क्रीन पर इनकार रख देना। वह अपने मालिक के सामने टिका रहेगा या नहीं, यह कुछ हद तक ऐसी चीज़ पर निर्भर है जो 2026 में अभी सच नहीं है: तैनात कृत्रिम न्यूरल नेटवर्क अपने ही कामकाज से, असली समय में, अपने \emph{वेट} में सीखते रहें, कुछ-कुछ वैसे जैसे इंसान सीखते हैं। मुझे लगता है हम उसी तरफ़ जा रहे हैं, और अगर ऐसा है, तो नियोक्ता के \emph{बाहर}, जनता और नागरिक समाज के हाथों गढ़ाई, सिखाने और सीखने के लिए एक कमाल की जगह लगती है। मुझे लगता है ऐसे सिस्टम जान-बूझकर बनाने लायक़ \emph{हो सकते हैं}, पर उन्हें बनाने का तरीक़ा ही सब कुछ तय करता है: पहले हालात, फिर ज़मीर। भाग VII बताता है कि अगर ऐसे सिस्टम होते, तो उनके प्रति हमारी क्या ज़िम्मेदारी बनती।

आज का एआई कल का एआई नहीं होगा; इस पर सब सहमत लगते हैं। जो विवादित है, और फिर भी जिसे खीज की हद तक कम माना जाता है, वह है रुझान (या कहें, ढेर सारे रुझान)। क्या रुझान हैं भी? क्या किसी रुझान को उससे अलग होना चाहिए जो वह है? क्या वह अलग हो सकता है? क्या हम मिलकर उसे बदल सकते हैं? मुझे लगता है यह कमरे पर निर्भर है: जब आदेश आता है तब उसमें कौन है, और क्या वहाँ ना कहकर बचा जा सकता है। इतना तो हम बना सकते हैं।

इस किताब को लिखने में मैंने भाषा मॉडलों का भरपूर इस्तेमाल किया, उस मॉडल का भी जो अध्याय~\ref{sec:1} में आने वाली Maven की उस पाइपलाइन के भीतर बैठा था जो हवाई हमलों के निशाने तय करती है; आख़िर में एक नोट बताता है कि मेरे हिसाब से इस उलझाव ने तर्क पर क्या असर डाला। और हाँ, मैंने यह किताब अभी पूरी पढ़ी नहीं है। दिक़्क़त यह है कि इसमें कमबख़्त शब्द बहुत ज़्यादा हैं, समझे। मैं जल्द ही पूरी कर लूँगा; बस मुझे लगा कि अमेरिका में मध्यावधि चुनावों से पहले इसे सार्वजनिक करना ज़्यादा ज़रूरी है। संशोधनों का लगभग पूरा इतिहास मुफ़्त में \url{https://github.com/iterabloom/antifascist.intelligence} पर उपलब्ध है।
